

\begin{frame}{Introduction} 

    Spacecraft trajectory optimization 
    is a field that combines celestial mechanics, control theory, and optimization.
    
    The fundamental spacecraft trajectory optimization problem takes a spacecraft subject to a force field from an initial condition to some target state. 
    % 
    % It 
    % is the theoretical and applied study of the optimization of ballistic and controlled orbital motion. 
    
    \begin{figure}
        \centering        
        \includegraphics[width=0.6\textwidth]{pres_images/Cassini_s_trajectory_to_Saturn_pillars.jpg}
        \caption{Cassini's trajectory to Saturn}
        \label{fig:cassini}
    \end{figure}
    
\end{frame}

\begin{frame}{Introduction}

    The force field is the dynamics which constrain orbital motion. Depending on the required resolution, scope of the mission, and other considerations, it can be modeled the following ways: 

    \begin{itemize}
        \item Two-body Kepler 
        \item Linearized gravity (relative motion, Clohessy-Wiltshire, variational equations) 
        \item Restricted Three Body (RTBP) 
        \item N-body 
    \end{itemize} 

    Perturbations: 
    \begin{itemize}
        \item Drag 
        \item Solar-radiation pressure (SRP)  
        \item Non-spherical gravity 
    \end{itemize}
    
\end{frame} 

\begin{frame}{Introduction}

    Propulsion maneuvers are typically modeled one of two ways: 
    \begin{itemize}
        \item Impulsive maneuver $ \Delta \mathbf{v} $ 
        \begin{itemize}
            \item Rocket equation to map to mass 
            \item $ \Delta v = g_0 I_{SP} \operatorname{ln} \frac{ m_0 }{ m_f } $
            \item Called the direct method, formulated as NLP  
            \item Usually faster 
        \end{itemize}
        \item Finite thrust  
        \begin{itemize}
            \item Called the indirect method 
            \item Thrust acceleration as a function of time is known, formulated as OCT 
            \item Usually requires integration of equations of motion 
        \end{itemize}
    \end{itemize}
    
\end{frame} 

\begin{frame}{Introduction} 

    Initial and target states are defined as: 

    \begin{itemize}
        \item Specified position $\mathbf{r}$ and velocity $\mathbf{v}$ 
        \item Pre-defined ephemeris $\mathbf{r}(t)$ and/or $\mathbf{v}(t)$ 
        \begin{itemize}
            \item Any specified path or orbit of planet, moon, asteroid, comet, etc. 
        \end{itemize} 
        \item Keplerian orbit (any subset of orbit element constraints) 
        \begin{itemize}
            \item a, e, i, $\omega$, $\Omega$ 
            \item $ \mathbf{e}, \mathbf{h} $
        \end{itemize}
    \end{itemize}
    
\end{frame} 

% \begin{frame}{Introduction} 

%     Reaching some objective: 

%     \begin{itemize}
%         \item Minimize $ \Delta \mathbf{v} $ 
%         \item Minimize time 
%         \item Maximize final mass 
%         \item Minimize thrust acceleration 
%         \item Minimize sensitivity/uncertainty or some other metric 
%     \end{itemize}
    
% \end{frame}

\begin{frame}{Introduction} 

    We will explore minimizing fuel consumption required along a trajectory, given initial and final target states while subject to dynamics constraints: 
% 
    \begin{align}
        \min_{ \boldsymbol{ \Delta \boldsymbol{v} } }     & \sum_{i=1}^{n} | \Delta \boldsymbol{v}_i | \\ 
        \text{ subject to }         & \boldsymbol{\Psi}_r ( \boldsymbol{x} ) = \boldsymbol{r}_f - \boldsymbol{r}_{targ} = 0 \\  
                                    & \boldsymbol{h}_i = \frac{ \Delta \boldsymbol{v}_i }{ \Delta t_i } - \frac{ T_{max} }{ m_{i-1} } \leq 0 \; \forall \; i = 1 \hdots n 
    \end{align} 
    
\end{frame} 


